\documentclass[aspectratio=169,10pt]{beamer}
\usetheme[numbering=fraction,progressbar=frametitle]{metropolis}
\usepackage{fontspec}
\setsansfont{Latin Modern Sans}
\usepackage{amsmath,amssymb,mathtools,graphicx,booktabs}
\usepackage{polyglossia}
\setdefaultlanguage{french}
\setbeamertemplate{navigation symbols}{}
\setbeamerfont{frametitle}{size=\large}
\setbeamerfont{normal text}{size=\normalsize}
\setlength{\parskip}{.45em}
\hypersetup{pdfauthor={},pdfsubject={Lecture slides},pdfkeywords={},pdftitle={Droites du plan}}
\providecommand{\N}{\mathbb{N}}
\providecommand{\Z}{\mathbb{Z}}
\providecommand{\D}{\mathbb{D}}
\providecommand{\Q}{\mathbb{Q}}
\providecommand{\R}{\mathbb{R}}
\providecommand{\C}{\mathbb{C}}
\providecommand{\vect}[1]{\overrightarrow{#1}}
\providecommand{\vi}{\vec{\imath}}
\providecommand{\vj}{\vec{\jmath}}
\providecommand{\vk}{\vec{k}}
\providecommand{\norm}[1]{\left\lVert #1\right\rVert}
\providecommand{\vc}[2]{\begin{pmatrix}#1\\#2\end{pmatrix}}
\providecommand{\deter}[2]{\det\left(#1,#2\right)}
\providecommand{\abs}[1]{\left\lvert #1\right\rvert}
\providecommand{\dd}{\,\mathrm{d}}
\providecommand{\e}{\mathrm{e}}
\providecommand{\Card}{\operatorname{Card}}
\providecommand{\E}{\operatorname{E}}
\providecommand{\V}{\operatorname{V}}
\providecommand{\degre}{^{\circ}}
\providecommand{\un}[1]{\,\mathrm{#1}}
\providecommand{\sci}[2]{#1\times10^{#2}}
\providecommand{\pgcd}{\operatorname{PGCD}}
\providecommand{\ppcm}{\operatorname{PPCM}}
\title{Droites du plan}
\subtitle{Cours et exemples}
\date{}
\author{}
\begin{document}
\begin{frame}[plain]
\vfill
{\usebeamerfont{title}\inserttitle\par}
\medskip
{\usebeamerfont{subtitle}\insertsubtitle\par}
\medskip
{\usebeamercolor[fg]{progress bar}\rule{\textwidth}{.5pt}}
\vfill
\end{frame}
\begin{frame}{Vecteur directeur et équation cartésienne}
Un vecteur directeur d'une droite est un vecteur \textbf{non nul} ayant la direction de cette droite. Si \(A\ne B\), \(\vect{AB}\) dirige \((AB)\).
Toute droite admet une équation \(ax+by+c=0\) avec \((a;b)\ne(0;0)\). Un vecteur directeur est \(\vc{-b}{a}\).
\end{frame}
\begin{frame}{Vecteur directeur et équation cartésienne}
Pour \(A(-2;4)\) et \(B(6;2)\), calculer \(\vect{AB}\) puis vérifier que \(\vec u(4;-1)\) dirige \((AB)\).
\end{frame}
\begin{frame}{Correction : Vecteur directeur et équation cartésienne}
\(\vect{AB}=(8;-2)=2\vec u\). Le vecteur \(\vec u\) est non nul et colinéaire à \(\vect{AB}\) ; il dirige donc \((AB)\).
\end{frame}
\begin{frame}{Vecteur directeur et équation cartésienne}
La droite passant par \(A(x_A;y_A)\) et dirigée par \(\vec u(\alpha;\beta)\ne\vec0\) a pour équation
\[\beta(x-x_A)-\alpha(y-y_A)=0.\]
\end{frame}
\begin{frame}{Vecteur directeur et équation cartésienne}
Un point \(M(x;y)\) appartient à la droite si et seulement si \(\vect{AM}\) est colinéaire à \(\vec u\), y compris pour \(M=A\). Or
\[\det(\vect{AM},\vec u)=(x-x_A)\beta-(y-y_A)\alpha.\]
L'annulation du déterminant donne l'équation annoncée. En développant, \(a=\beta\), \(b=-\alpha\) et \(c=-\beta x_A+\alpha y_A\) ; comme \(\vec u\ne\vec0\), \((a;b)\ne(0;0)\).
\end{frame}
\begin{frame}{Vecteur directeur et équation cartésienne}
Le vecteur nul ne dirige aucune droite. Le signe du déterminant dépend de l'ordre des vecteurs, mais son annulation n'en dépend pas.
\end{frame}
\begin{frame}{Déterminer une équation de droite}
À partir d'un point et d'un vecteur directeur : choisir \(a,b\) puis remplacer \(x,y\) par les coordonnées du point pour trouver \(c\). À partir de deux points distincts : calculer d'abord un vecteur directeur.
\end{frame}
\begin{frame}{Déterminer une équation de droite}
Déterminer une équation cartésienne de la droite passant par \(A(2;-1)\) et dirigée par \(\vec u(3;2)\).
\end{frame}
\begin{frame}{Correction : Déterminer une équation de droite}
\(a=2\), \(b=-3\) conviennent. La droite a donc une équation \(2x-3y+c=0\). Comme \(A\) lui appartient, \(4+3+c=0\), d'où \(c=-7\). Une équation est \(2x-3y-7=0\). Vérification : \(2\times2-3\times(-1)-7=0\) et \((3;2)\) est bien directeur.
\end{frame}
\begin{frame}{Déterminer une équation de droite}
Déterminer une équation cartésienne de \((AB)\) pour \(A(3;3)\) et \(B(2;-2)\).
\end{frame}
\begin{frame}{Correction : Déterminer une équation de droite}
\(\vect{AB}=(-1;-5)\) dirige la droite. On peut choisir \((a;b)=(5;-1)\). Avec \(A\), \(15-3+c=0\), donc \(c=-12\). Ainsi \(5x-y-12=0\). Pour \(B\), \(10+2-12=0\), ce qui confirme le résultat.
\end{frame}
\begin{frame}{Déterminer une équation de droite}
Le point \(P(1;2)\) appartient-il à \(d:2x-3y+4=0\) ? Le point \(Q(2;1)\) lui appartient-il ?
\end{frame}
\begin{frame}{Correction : Déterminer une équation de droite}
Pour \(P\), \(2-6+4=0\) : oui. Pour \(Q\), \(4-3+4=5\ne0\) : non.
\end{frame}
\begin{frame}{Équation réduite, pente et fonction affine}
Une droite non verticale a une unique équation réduite \(y=mx+p\). \(m\) est la pente (coefficient directeur), \(p\) l'ordonnée à l'origine. Elle représente la fonction affine \(x\mapsto mx+p\) ; \((1;m)\) est un vecteur directeur.
Une droite verticale a une équation \(x=k\) ; elle n'est pas la courbe d'une fonction de \(x\). Une droite horizontale a une équation \(y=k\).
\end{frame}
\begin{frame}{Équation réduite, pente et fonction affine}
Si \(x_A\ne x_B\), la pente de \((AB)\) vaut \(m=\dfrac{y_B-y_A}{x_B-x_A}\). Pour trouver \(p\), on utilise ensuite \(y_A=mx_A+p\).
\end{frame}
\begin{frame}{Équation réduite, pente et fonction affine}
\begin{columns}[T,onlytextwidth]
\begin{column}{.57\textwidth}
Déterminer l'équation réduite de \(2x+3y-1=0\), puis tracer \(d:y=-2x+3\).
\end{column}
\begin{column}{.4\textwidth}
\centering
\includegraphics[width=\linewidth,height=4.5cm,keepaspectratio]{figures/seconde-03-1.pdf}
\end{column}
\end{columns}
\end{frame}
\begin{frame}{Correction : Équation réduite, pente et fonction affine}
\(3y=-2x+1\), donc \(y=-\dfrac23x+\dfrac13\). Pour \(d\), placer \((0;3)\) et \((1;1)\), puis tracer la droite ; un déplacement de \(1\) en abscisse correspond à \(-2\) en ordonnée.
\end{frame}
\begin{frame}{Équation réduite, pente et fonction affine}
Trouver la fonction affine \(f\) telle que \(f(0)=3\) et \(f(-1)=5\), puis la fonction affine \(g\) dont la courbe passe par \(C(-1;2)\) et \(E(5;-2)\).
\end{frame}
\begin{frame}{Correction : Équation réduite, pente et fonction affine}
\(f(x)=-2x+3\), car \(m=(3-5)/(0+1)=-2\). Pour \(g\), \(m=(-2-2)/(5+1)=-2/3\) et \(p=2-2/3=4/3\), donc \(g(x)=-\dfrac23x+\dfrac43\).
\end{frame}
\begin{frame}{Positions relatives de deux droites}
Deux droites sont sécantes, strictement parallèles ou confondues. Pour \(y=mx+p\) et \(y=m'x+p'\) :

\begin{itemize}
\item si \(m\ne m'\), elles sont sécantes ;
\item si \(m=m'\) et \(p\ne p'\), elles sont strictement parallèles ;
\item si \(m=m'\) et \(p=p'\), elles sont confondues.

\end{itemize}

Pour des équations cartésiennes, on compare les directions avec le déterminant. Si les directions sont colinéaires, on teste ensuite l'appartenance d'un point pour distinguer droites confondues et strictement parallèles.
\end{frame}
\begin{frame}{Positions relatives de deux droites}
Déterminer une équation de \(d\) passant par \(A(2;3)\) et parallèle à \(d':x-3y-3=0\).
\end{frame}
\begin{frame}{Correction : Positions relatives de deux droites}
\(d'\) a pour vecteur directeur \((3;1)\). On cherche \(d:x-3y+c=0\). Comme \(A\in d\), \(2-9+c=0\), donc \(c=7\). Ainsi \(d:x-3y+7=0\). Les deux droites sont strictement parallèles, car les constantes diffèrent pour les mêmes coefficients de \(x,y\).
\end{frame}
\begin{frame}{Positions relatives de deux droites}
Étudier les positions relatives :

\begin{enumerate}
\setcounter{enumi}{0}\item \(2x-y+1=0\) et \(4x-2y+2=0\) ;
\item \(x=2\) et \(x=-1\) ;
\item \(y=2x-3\) et \(y=-x+3\).

\end{enumerate}
\end{frame}
\begin{frame}{Correction : Positions relatives de deux droites}
1. La deuxième équation est le double de la première : droites confondues. 2. Deux verticales distinctes : strictement parallèles. 3. Les pentes \(2\) et \(-1\) diffèrent : sécantes. Leur intersection vérifie \(2x-3=-x+3\), donc \(x=2\) et \(y=1\).
\end{frame}
\begin{frame}{Intersection : résolution par substitution}
Résoudre un système de deux équations en \(x,y\), c'est trouver tous les couples \((x;y)\) vérifiant simultanément les deux équations. Pour deux équations de droites : des droites sécantes donnent un unique couple solution ; des droites strictement parallèles donnent aucune solution ; des droites confondues donnent une infinité de solutions.
\end{frame}
\begin{frame}{Intersection : résolution par substitution}
Isoler une inconnue dans une équation, remplacer son expression dans l'autre, résoudre puis retrouver la seconde inconnue. Vérifier le couple dans les deux équations initiales.
\end{frame}
\begin{frame}{Intersection : résolution par substitution}
Résoudre \(\begin{cases}-2x+y+3=0,\\3x-5y-2=0.\end{cases}\)
\end{frame}
\begin{frame}{Correction : Intersection : résolution par substitution}
La première équation donne \(y=2x-3\). En remplaçant dans la seconde : \(3x-5(2x-3)-2=0\), soit \(-7x+13=0\). Ainsi \(x=13/7\) puis \(y=5/7\). L'ensemble solution est \(S=\{(13/7;5/7)\}\). Les substitutions donnent \((-26+5+21)/7=0\) et \((39-25-14)/7=0\).
\end{frame}
\begin{frame}{Intersection : résolution par substitution}
Résoudre par substitution :
\[\text{(a) }\begin{cases}-2x+y+3=0\\x-4y+5=0\end{cases}\qquad
\text{(b) }\begin{cases}2x+3y-10=0\\4x+5y-16=0.\end{cases}\]
\end{frame}
\begin{frame}{Correction : Intersection : résolution par substitution}
(a) \(y=2x-3\) puis \(x-4(2x-3)+5=0\), d'où \(x=17/7\), \(y=13/7\).

(b) \(x=(10-3y)/2\) ; \(2(10-3y)+5y-16=0\), donc \(y=4\), puis \(x=-1\).
\end{frame}
\begin{frame}{Combinaison linéaire et Python}
Multiplier les équations par des nombres non nuls pour obtenir des coefficients opposés d'une inconnue ; les additionner pour éliminer cette inconnue. Conserver une équation du système pour calculer l'autre inconnue.
\end{frame}
\begin{frame}{Combinaison linéaire et Python}
Résoudre \(\begin{cases}2x+3y+1=0\\3x-4y-2=0.\end{cases}\) Reprendre ensuite les deux systèmes de l'exemple précédent par combinaison.
\end{frame}
\begin{frame}{Correction : Combinaison linéaire et Python}
Multiplier la première équation par \(3\) et la seconde par \(-2\) donne \(6x+9y+3=0\) et \(-6x+8y+4=0\). La somme donne \(17y+7=0\), d'où \(y=-7/17\). Puis \(2x-21/17+1=0\), donc \(x=2/17\).
Pour (a), ajouter quatre fois la première équation à la seconde : \(-7x+17=0\), donc \((x;y)=(17/7;13/7)\). Pour (b), soustraire deux fois la première à la seconde : \(-y+4=0\), donc \((x;y)=(-1;4)\).
\end{frame}
\begin{frame}[fragile]{Python}
\small
\begin{verbatim}
def determinant(u, v):
    return u[0]*v[1] - u[1]*v[0]

def alignes(A, B, C):
    AB = (B[0]-A[0], B[1]-A[1])
    AC = (C[0]-A[0], C[1]-A[1])
    return determinant(AB, AC) == 0
\end{verbatim}
\end{frame}
\begin{frame}{Combinaison linéaire et Python}
Prévoir \texttt{alignes((0,0), (2,1), (6,3))} puis \texttt{alignes((0,0), (2,1), (6,4))}.
\end{frame}
\begin{frame}{Correction : Combinaison linéaire et Python}
\texttt{True}, puis \texttt{False} : les déterminants sont \(0\) et \(2\). Le test exact convient ici aux coordonnées entières ; les approximations flottantes demandent de tenir compte des erreurs numériques.
\end{frame}
\end{document}
